\section{Inleiding}
Het hergebruik van overheidsgegevens veronderstelt dat afnemers niet alleen waarden kunnen ontvangen, maar ook hun betekenis en toepassingsvoorwaarden kunnen vaststellen. Het European Interoperability Framework (EIF) onderscheidt juridische, organisatorische, semantische en technische interoperabiliteit. De historische eisen aan Nederlandse basisregistraties verbinden registratie bovendien met onder meer wettelijke verankering, terugmelding, kwaliteitsborging en verantwoordelijkheden. Daarmee is gegevensuitwisseling een vraagstuk dat de technische interface overstijgt. \citep{S43,S25}

Voor afzonderlijke delen van dit vraagstuk bestaan uiteenlopende instrumenten. ISO 704 behandelt terminologisch werk, ISO/IEC 11179 metadata\-registratie, MIM informatiemodellering, SKOS begrippenstelsels en PROV-O provenance. Hun onderwerpen zijn verwant, maar hun representatie-eenheden verschillen. Zo is een begrip in SKOS niet zonder meer een OWL-klasse en is een data-element in ISO/IEC 11179 niet hetzelfde als een waarde in een record. \citep{S11,S08,S13,S31,S30,S33} De centrale ontwerpvraag is daarom welke expliciete verbindingen nodig zijn om deze instrumenten samenhangend toe te passen.

Deze paper stelt daarvoor een Semantic Master Data Product (SMDP) voor. Het begrip wordt hier als eigen werkdefinitie gebruikt: een bestuurlijk beheerde en versieerbare publicatie-eenheid die master- en referentiegegevens verbindt met machinevindbare betekenis, gezagsgrond, identiteit, historie, provenance, kwaliteitsregels en gebruiksmetadata. De productbenadering sluit aan bij de gebruikersoriëntatie van dataproductliteratuur; zij veronderstelt geen verplichte data-mesharchitectuur. \citep{L12} De voorgestelde integratie is het te onderzoeken artefact, niet een uit bestaande standaarden rechtstreeks voortvloeiend productschema.

AI wordt beschouwd als een aanvullende afnemer. FAIR motiveert machinegericht hergebruik, terwijl kennisgraafliteratuur technieken voor expliciete representatie en verwerking bespreekt. Geen van deze bronnen bewijst dat het hier voorgestelde product betere AI-antwoorden oplevert. \citep{L11,L13} Deze paper houdt die verwachting daarom gescheiden van de onderbouwde mogelijkheden van standaarden.

\section{Probleemstelling}
\subsection{Praktisch probleem}
Het praktische ontwerpprobleem is de mogelijke discontinuïteit tussen de betekenis die een bevoegde organisatie vaststelt en de betekenis die een afnemer uit een gegevenslevering afleidt. Een API kan een waarde structureel correct leveren zonder in dezelfde representatie de definitie, gezagsgrond, historische geldigheid of herkomst van die waarde beschikbaar te maken. OpenAPI beschrijft een interface; PROV-O beschrijft herkomstrelaties; NL-SBB beschrijft begrippen en hun beschrijving. Geen van deze toepassingsgebieden omvat zelfstandig de volledige keten. \citep{S19,S33,S14} Dit is een analyse van hun dekking, geen gemeten frequentie van fouten in het Stelsel.

Een illustratief risico is dat een afnemer een technisch veld \texttt{adres} als unieke objectidentiteit interpreteert, terwijl het veld in het aangeboden model slechts een locatiebeschrijving aanduidt. Dit voorbeeld is hypothetisch. Het laat zien waarom het ontwerp onderscheid moet kunnen uitdrukken tussen beschrijving, identifier, record en het bedoelde object. De theoretische motivatie ligt in terminologische onderscheidingen, metadata\-registratie en ontologische analyse van identiteit. \citep{S11,S08,L08,L10}

\subsection{Wetenschappelijke kennisvraag en antecedenten}
De brede claim dat masterdata nog niet met semantische technologie worden verbonden, is niet houdbaar. Wang et al. beschrijven al in 2009 SMDM met virtuele RDF-toegang, ontologieën en mappings. \citep{L18} DPROD biedt productmetadata en Gualo et al. operationaliseren masterdatakwaliteit. \citep{S51,L19} Deze tijdens de gerichte tegencontrole gevonden antecedenten maken semantische MDM, productbeschrijving en kwaliteitsmeting op zichzelf ongeschikt als nieuwheidsclaim.

De resterende kennisvraag luidt: wanneer levert expliciete, versiegebonden contextbinding van identiteit, tijd en gezagsgrond bij overheidsgegevens meer interpretatieve en beheermatige waarde dan dezelfde broninhoud in een eenvoudiger conventioneel contextregister? Het gaat om een vergelijkende en conditionele effectvraag. Het corpus toont niet aan dat het voorgestelde antwoord nog nergens bestaat. Een volledige review van verwante oplossingen blijft nodig; het praktische integratieprobleem mag niet worden verward met een bewezen wetenschappelijke lacune.

\begin{longtable}{@{}P{.22\linewidth}P{.32\linewidth}P{.38\linewidth}@{}}
\caption{Antecedenten en resterende vergelijkingsvraag; niet-vermelding van een functie is geen bewijs van afwezigheid.}\\
\toprule Antecedent & Bestaande bijdrage & Wat voor deze paper nog moet worden vastgesteld \\ \midrule
SMDM \citep{L18} & Semantische toegang tot masterdata en koppeling met ontologieën. & Of extra contextbinding bij tijd/gezag boven een bestaande semantische MDM-benadering nodig en nuttig is. \\
MDQM \citep{L05}; kwaliteitsmodel \citep{L19} & Organisatorische inrichting en systeemondersteuning en masterdatakwaliteitsmetingen. & Welke consumentuitkomsten boven correct kwaliteitsbeheer veranderen. \\
DPROD \citep{S51} & Productbeschrijving met doel, eigenaar, lifecycle en verbindingen met datasets/diensten. & Welk minimaal aanvullend profiel nodig is voor de gekozen uitspraakcontext; geen generiek productschema opnieuw uitvinden. \\
Governance en kennisgrafen \citep{L07,L13} & Bestaande governanceconcepten en technieken voor context en verwerking. & Welke combinatie meerwaarde heeft bij expliciet gemeten kosten en rivaliserende implementaties. \\
\bottomrule\end{longtable}

De kandidaatbijdrage is dus geen nieuwe opslagtechniek of universele ontologie. Zij bestaat uit een expliciet contextbindingscontract, afgebakende mechanismen en een experiment dat ook een eenvoudiger ontwerp kan laten winnen. In termen van DSR blijft nog te bepalen of dit verbetering, contextuele toepassing of routineontwerp is. \citep{L04}

\subsection{Maatschappelijke en wetenschappelijke relevantie}
De maatschappelijke relevantie ligt in verantwoord hergebruik van gegevens waarvoor publieke verantwoordelijkheden en kwaliteitsafspraken bestaan. De eisen aan basisregistraties en de kwaliteitsbaseline maken terugmelding en kwaliteitszorg tot expliciete aandachtspunten. \citep{S25,S26} Een ontwerp dat zichtbaar maakt welke uitspraak door wie en op welke grond wordt gedragen, kan deze verantwoordelijkheden mogelijk ondersteunen. Minder fouten, lagere kosten of beter overheidsoptreden worden hier niet als gerealiseerde effecten geclaimd.

De wetenschappelijke relevantie ligt in het expliciteren en falsifieerbaar maken van integratiebeslissingen. Het voorgestelde metamodel maakt onderscheidingen toetsbaar die in afzonderlijke standaarden verschillende vormen hebben. Een toekomstige evaluatie kan vervolgens onderzoeken welke verbindingen bijdragen aan interpretatiekwaliteit en tegen welke beheerkosten. Daarmee richt de bijdrage zich op overdraagbare ontwerpprincipes, naast de specifieke implementatie. \citep{L01,L03,L04}

\section{Onderzoeksdoel en onderzoeksvragen}
\label{sec:questions}
Het doel is een onderbouwd, falsifieerbaar ontwerpvoorstel te formuleren. De centrale these T luidt: bij taken die identiteit, tijd en gezag combineren, kan expliciete contextbinding netto meerwaarde bieden boven dezelfde informatie in een conventioneel register. Noodzaak van RDF, centrale opslag, alle vijftien voorzieningen of OntoUML maakt geen deel uit van T. Uit de alternatieve onderzoekslijnen in het repositorydossier wordt de integratiegerichte lijn gekozen; ontologische verdieping wordt als ontwerpoptie onderzocht en governance als doorsnijdende verantwoordelijkheid behandeld.

\noindent\textbf{Hoofdvraag.} Hoe kan een Semantic Master Data Product voor de Nederlandse overheid worden ontworpen waarin nationale, Europese en internationale standaarden betekenis, gezag, identiteit en historie traceerbaar verbinden met masterdata en machine-interpreteerbare distributie, en hoe kan de meerwaarde daarvan worden geëvalueerd binnen het Stelsel van Basisregistraties?

\begin{description}[style=nextline,leftmargin=0pt,labelwidth=0pt]
\item[OV1 --- Eisen en toepasselijkheid.] Welke inhoudelijke eisen en toepasselijkheidsvoorwaarden volgen uit het geselecteerde corpus van literatuur, standaarden en overheidskaders?
\item[OV2 --- Abstractieniveaus en semantische hiaten.] Welke onderscheidingen, overlappingen en conflicterende interpretaties moeten worden gemodelleerd om onterechte equivalentie te voorkomen?
\item[OV3 --- Integratie.] Welke getypeerde relaties verbinden definities, modellen, data-elementen, identiteit, tijd, provenance en epistemische status over versies heen?
\item[OV4 --- Ontwerp en beheer.] Welke architectuur, ontwerpprincipes en governanceafspraken kunnen deze relaties operationaliseren, en welke alternatieven zijn verdedigbaar?
\item[OV5 --- Toetsbare meerwaarde.] Hoe kunnen interpretatie, consistentie, traceerbaarheid en beheerbaarheid worden vergeleken met technische masterdata zonder expliciete semantische verbindingen?
\item[OV6 --- Machineconsumptie.] Onder welke voorwaarden kunnen informatiesystemen en AI de betekenis en bewijsstatus correct benutten, en hoe moet dat worden getoetst?
\end{description}
De paper beantwoordt OV1--OV4 met een afgebakende synthese en ontwerpvoorstel. Voor OV5 en OV6 levert zij een evaluatieontwerp; effectvragen blijven open.

\section{Onderzoeksmethode}
\subsection{Design Science Research en onderzoeksstatus}
Design Science Research (DSR) verbindt de ontwikkeling van artefacten met evaluatie en kennisontwikkeling. De methode van Peffers et al. ordent probleemidentificatie, oplossingsdoelen, ontwerp, demonstratie, evaluatie en communicatie. \citep{L01,L02} Deze paper rapporteert probleemafbakening, doelen en een voorgesteld ontwerp. Het BAG-fragment is een analytische illustratie van modelleringsvragen, geen operationele demonstratie. De formele empirische evaluatie is nog niet uitgevoerd. Het bestaan van een prototype maakt deze fase niet tot een afgeronde DSR-cyclus.

Het DSR-artefact bestaat uit vier samenhangende onderdelen: een getypeerd metamodel, een reference architecture, negen ontwerpprincipes en een evaluatieprotocol met traceability. Een implementatie kan deze onderdelen later operationaliseren. Voor de evaluatie wordt FEDS gebruikt om formatieve en summatieve, kunstmatige en praktijkgerichte episodes te onderscheiden. \citep{L03} De uitwerking van episodes, taken en metrieken is een eigen voorstel.

\subsection{Bronbasis en syntheseprocedure}
Het onderzoek vertrok uit het repositorydossier met 50 standaard-/kadervermeldingen en 17 literatuurvermeldingen. Een afzonderlijk vastgelegde kritische review voegde één productontologie en twee academische antecedenten toe, vóór herziening van de argumentatie: S51, L18 en L19. De herziene matrices bevatten 51 respectievelijk 19 vermeldingen; dit zijn geen 70 onafhankelijke studies. De zoekstrings en selectie zijn vastgelegd in \texttt{review/search-log.md}. De afbakening is doelgericht; er is geen reproduceerbare databasebrede selectie beschikbaar waarmee volledigheid kan worden geclaimd.

De synthese verloopt van geregistreerde bronclaim naar ontwerpprobleem, requirement, principe, architectuurverantwoordelijkheid en toekomstige metriek. Daarbij wordt eerst vastgesteld op welke eenheid een bron betrekking heeft, bijvoorbeeld begrip, data-element of dienst. Vervolgens wordt onderzocht welke relatie met de aangrenzende eenheid ontbreekt. Pas daarna wordt een integratiekeuze voorgesteld. De tabellen in de bijlage maken deze redenering controleerbaar, maar vervangen geen inhoudelijke deskundigenbeoordeling.

Bronnen worden naar bewijssterkte behandeld. Volledige of relevante primaire tekst ondersteunt alleen de daadwerkelijk vastgelegde inhoud. Abstracts ondersteunen alleen bevindingen op abstractniveau. Officiële normcatalogi ondersteunen scope- en editieclaims, niet clausuleconformiteit. L06 is uitsluitend bibliografisch geverifieerd en wordt niet als inhoudelijk bewijs gebruikt. Het onderliggende corpus bevat bovendien conceptuele bijdragen, empirische studies en normatieve teksten; een normatieve aanbeveling wordt niet omgezet in een empirische effectclaim.

\subsection{Epistemische scheiding en onderzoeksintegriteit}
Zes categorieën blijven onderscheiden. \status{Bronbevindingen} betreffen wat het dossier over literatuur en standaarden vastlegt. \status{Requirements} zijn daaruit afgeleide ontwerpverplichtingen. \status{Ontwerpprincipes} generaliseren die verplichtingen. De \status{reference architecture} is een voorgestelde realisatie op logisch niveau. Het \status{prototype} is een bestaand onderzoeksartefact waarvan in dit dossier geen implementatiedetails zijn aangetoond. De \status{empirische evaluatie} moet nog plaatsvinden.

De paper vult het register aan met R23 voor expliciete epistemische metadata, DP1--DP9 voor ontwerpprincipes en P01--P06 voor ontwerp- en afbakeningsproposities. De latere review introduceert drie expliciet geregistreerde externe antecedenten en TP1--TP4 als niet-getoetste propositions. Zij is door dezelfde AI-assistent uitgevoerd, niet door een organisatorisch onafhankelijke reviewer. Het oorspronkelijke manuscript en de volledige review zijn vóór de revisie bewaard. De identifiers R01--R23 blijven behouden. Een afzonderlijk \texttt{derivation-register.csv} maakt aanvullende aannames, alternatieven en experimentkoppelingen expliciet. Een redactionele audit toetst uitsluitend of deze afleidingsketen intern gesloten is; zij valideert niet de inhoudelijke juistheid van het voorgestelde ontwerp.

\subsection{Afleiding en DSR-beslismomenten}
De afleiding onderscheidt vier beweringstypen: een bron beschrijft een constructie; een norm stelt binnen scope een eis; de ontwerper kiest een aanvullende contractregel; een hypothese voorspelt een effect. De overgang van de eerste twee naar de laatste twee vereist een expliciete aanname. Een beschikbare provenanceproperty bijvoorbeeld maakt context representabel, maar verplicht geen uitspraakgebonden register en voorspelt geen betere antwoorden. \citep{S33}

DSR wordt gebruikt als onderzoeksorganisatie, niet als kwaliteitsstempel. De uitgevoerde activiteiten zijn dossieranalyse, brontegencontrole en specificatie van een voorstel. De analytische voorbeelden hieronder zijn geen demonstratie van software. Gate G-A verlangt bron- en constructreview; G-B een bevroren minimale implementatie en geadjudiceerde taakset; G-C een preregistratie van uitkomsten, marges en analyse; G-D uitvoering van vergelijkende en praktijkgerichte evaluatie. Geen van deze gates wordt door compilatie of administratieve traceability vervuld. \citep{L02,L03}
